% ============================================================
% Capitolo 5 — Formalismi a stati finiti, seconda parte
% (20260410 Formalismi a stati finiti 2p.pdf)
% ============================================================
\chapter{Formalismi a stati finiti (seconda parte): Part-Whole Statecharts}
\label{cap:formalismi-2}

\section{Part-Whole Statecharts (PWS)}
\begin{defn}[Part-Whole Statecharts]
Formalismo che \textbf{separa il comportamento dei componenti da quello del sistema complessivo}, attraverso due sezioni distinte e parallele:
\begin{enumerate}
    \item \textbf{Assembly} (comportamenti locali dei componenti);
    \item \textbf{System section} (\emph{Whole}, comportamento globale).
\end{enumerate}
\end{defn}

\subsection{Observer/Controller Pattern}
Il pattern di riferimento dei PWS risponde alla domanda: «chi può conoscere/controllare gli altri componenti?»
\begin{itemize}
    \item \textbf{Componenti}: no, neppure il Sistema (Whole) è conosciuto dai componenti. Ogni componente è \textbf{comportamentalmente isolato};
    \item \textbf{Sistema}: sì, conosce e controlla i componenti.
\end{itemize}

\section{Come funzionano i PWS}
Ogni componente è modellato come un'\textbf{automa a stati riusabile e testabile in isolamento}. Il sistema (\emph{Whole}) coordina i componenti tramite:
\begin{enumerate}
    \item \textbf{Eventi diretti} (non broadcast);
    \item \textbf{Guardie} sulle transizioni.
\end{enumerate}

\subsection{Componenti: attuatori e sensori}
\begin{itemize}
    \item \textbf{Attuatori}: caratterizzati da eventi \textbf{trigger} (come \texttt{start} e \texttt{stop}) che fanno scattare le rispettive transizioni;
    \item \textbf{Sensori}: possono avere \textbf{transizioni che scattano da sole}, senza alcun trigger che le faccia scattare (\textbf{transizioni autonome}, notate con notazione grafica dedicata).
\end{itemize}

\subsection{Sistema (intero, \emph{Whole})}
\begin{itemize}
    \item È caratterizzato da eventi trigger (come \texttt{stop} e \texttt{on}) che fanno scattare le rispettive transizioni;
    \item può avere \textbf{guardie} alle transizioni, che costituiscono condizione necessaria affinché la transizione scatti;
    \item può inviare \textbf{eventi diretti} (azioni) alle componenti attraverso \textbf{liste di azioni}: ad esempio \texttt{< p.stop, v.close >}.
\end{itemize}

\section{Proposizioni di stato e verifica della sicurezza}
Attraverso \textbf{proposizioni di stato} si possono:
\begin{itemize}
    \item esplicitare \textbf{invarianti} (es.\ «se pressione è alta, la pompa è spenta e la valvola è chiusa»);
    \item derivare \textbf{safety axioms} controllando solo gli stati del sistema (non combinazioni esplosive di stati locali);
    \item evitare il \textbf{model checking pesante}, usando logica proposizionale per garantire sicurezza.
\end{itemize}

\section{Statecharts vs.\ PW-Statecharts}
\begin{note}[Confronto]
\begin{tabular}{p{0.45\textwidth}p{0.45\textwidth}}
\textbf{Problema Statecharts} & \textbf{Soluzione PW-Statecharts}\\
\hline
Perdita di modularità & Isolando completamente i componenti\\
Complessità della sincronizzazione & Centralizzando logica e coordinamento\\
Difficoltà di test e certificazione & Permettendo test modulari e semantica globale chiara\\
Verifica della sicurezza & Tramite proposizioni e transizioni sicure\\
\end{tabular}
\end{note}

\begin{intuition}[In sintesi]
Dove gli Statecharts classici di Harel fallivano (regioni ortogonali interdipendenti, Capitolo~\ref{cap:formalismi-1}), i PWS ribaltano l'architettura: i componenti sono automi completamente isolati e riusabili, e tutte le interazioni passano dal sistema (Whole), che osserva e comanda. La sicurezza si verifica a livello di sistema con proposizioni di stato, senza esplorare il prodotto cartesiano degli stati locali.
\end{intuition}
